\chapter*{术语对照（译本附加）}
\addcontentsline{toc}{chapter}{术语对照（译本附加）}
本表汇集译本采用的主要术语。具体含义以各篇作者的定义为准。
\begin{longtable}{@{}p{.43\textwidth}p{.53\textwidth}@{}}
\toprule 中文 & English \\\midrule\endfirsthead
\toprule 中文 & English \\\midrule\endhead
\bottomrule\endfoot
约化群 & reductive group \\
半单群 & semisimple group \\
环面 & torus \\
分裂／拟分裂 & split / quasi-split \\
根数据／根系 & root datum / root system \\
余根 & coroot \\
Weyl 群 & Weyl group \\
Borel 子群 & Borel subgroup \\
抛物子群 & parabolic subgroup \\
幺幂 & unipotent \\
根基 & radical \\
特征／拟特征 & character / quasi-character \\
容许表示 & admissible representation \\
光滑表示 & smooth representation \\
缓增表示 & tempered representation \\
逆反表示 & contragredient representation \\
交织算子 & intertwining operator \\
尖点／超尖点／强尖点 & cuspidal / supercuspidal / hypercuspidal \\
自守形式 & automorphic form \\
自守表示 & automorphic representation \\
阿代尔／伊代尔 & adele / idele \\
Weil 群 & Weil group \\
惯性群 & inertia group \\
Frobenius 元素 & Frobenius element \\
互反律 & reciprocity law \\
基变换 & base change \\
轨道积分 & orbital integral \\
迹公式 & trace formula \\
振子表示 & oscillator representation \\
亚辛群 & metaplectic group \\
约化对偶对 & dual reductive pair \\
Shimura 簇 & Shimura variety \\
模簇 & moduli variety \\
典范模型 & canonical model \\
动机 & motive \\
上同调 & cohomology \\
平展 & étale \\
局部域／整体域 & local field / global field \\
数域／函数域 & number field / function field \\
非分歧 & unramified \\
\end{longtable}
